\chapter{Conclusion}

\section{Overview}

The dissertation examined how prompt engineering affects correction behaviour in large language model (LLM)-based grammatical error correction (GEC). I focused on two questions at the same time: whether a prompt improves ERRANT accuracy, and whether it changes the learner sentence more than the correction task appears to require.

To study this, I combined ERRANT-based grammatical metrics with OCI. This allowed the baseline, instruction, role, and few-shot prompts to be compared both as correction methods and as different editing behaviours.

\section{Summary of Findings}

The results point to a clear pattern: prompt wording affected both grammatical correction performance and indicators of over-correction risk in this experimental setup.

Prompt refinement showed that small wording changes can alter the model outputs. The selected structured prompts performed better than the simple baseline in the later comparison. This suggests that the unconstrained prompt gave the model too much freedom when deciding how far to rewrite.

The final strategy comparison showed that the prompts did not fail or succeed in the same way. The baseline prompt produced the lowest $F_{0.5}$ score and the highest OCI, which is consistent with frequent but poorly targeted editing. The instruction prompt reduced OCI while maintaining strong performance. The role prompt gave the best balance in this study, with high precision, strong $F_{0.5}$, low OCI, and the lowest average edit distance. Few-shot prompting achieved the highest $F_{0.5}$, but its higher OCI suggests a broader correction scope.

Sentence length also affected the results. Across prompt conditions, long sentences had lower $F_{0.5}$ and higher OCI than short or medium sentences. This suggests that over-correction risk is affected by input complexity as well as by prompt design.

The robustness analysis showed that the broad OCI ranking did not depend on one weighting scheme. The baseline remained the highest-risk condition, while role and instruction prompting stayed near the lowest OCI values.

The trade-off analysis was useful because higher accuracy did not always require higher OCI. The role prompt achieved strong correction performance while keeping OCI low. In contrast, the baseline had both weak $F_{0.5}$ and the highest OCI, so more editing did not automatically mean better correction.

The qualitative examples made the numerical pattern easier to interpret. The baseline sometimes paraphrased, changed style, or produced task-format failures. These examples help explain why ERRANT alone was not enough for this project.

\section{Contributions}

The dissertation makes three main contributions.

First, it provides empirical evidence, from the W\&I + LOCNESS subset used here, that prompt wording can change LLM-based GEC behaviour. The comparison is limited to one fixed OpenAI GPT-4o configuration, but the differences between baseline, instruction, role, and few-shot prompts are visible across ERRANT, OCI, and edit-distance results.

Second, it applies OCI as an indicator of over-correction risk. The metric estimates source-output divergence while reducing the penalty for changes that preserve meaning and improve fluency. In practice, OCI made it possible to discuss rewriting behaviour alongside ERRANT accuracy instead of treating $F_{0.5}$ as the only outcome.

Third, it gives a practical interpretation of the prompt strategies tested in the project. The baseline behaved as the least controlled condition, the instruction prompt was more conservative, the role prompt was the strongest balanced condition, and few-shot prompting was the most accuracy-oriented but slightly more interventionist.

\section{Limitations}

The study has several limitations.

The experiments used a single dataset, W\&I + LOCNESS. Although this dataset is widely used in GEC research, the findings may not fully generalise to other learner corpora or domains.

The study also used a single LLM configuration: OpenAI GPT-4o with fixed decoding settings. Other models may respond differently to the same prompts, especially if they have different instruction tuning or decoding behaviour.

OCI remains an indirect measure rather than an explicit judgement of whether individual edits are unnecessary. Although the formulation accounts for fluency gain and meaning preservation, it cannot definitively determine whether individual edits are necessary, nor can it fully distinguish beneficial stylistic improvements from undesirable rewriting in all cases.

The study was also conducted in a controlled evaluation setting. A real GEC tool would involve user interaction, feedback presentation, and more varied input quality.

\section{Future Work}

Several follow-up studies would strengthen this work.

One useful next step would be to test the same prompt families across multiple models and datasets. This would show whether the role and instruction prompts remain strong outside the model and W\&I + LOCNESS subset used here.

Another direction is to develop more fine-grained over-correction metrics. OCI is useful for comparing prompt conditions, but it cannot decide on its own whether a specific stylistic change is pedagogically helpful. Semantic or discourse-level features may make this distinction clearer.

Future work could also test adaptive prompts that change according to sentence length or complexity. The sentence-length experiment suggests that longer inputs may need stricter preservation instructions.

User feedback would also make the evaluation more realistic. Learners or teachers could judge whether a correction is helpful, too minimal, or too interventionist, which automatic metrics cannot fully determine.

\section{Final Remarks}

The main finding is that prompt engineering affects more than output accuracy. In these experiments, prompts also changed how much the model rewrote the learner sentence. Evaluating both ERRANT performance and OCI therefore gave a clearer picture of correction behaviour than accuracy alone.

For learner-facing GEC, a good prompt should not simply encourage more edits. It should help the model correct grammatical errors while preserving as much of the learner's expression as possible.
